\chapter{Methods and Experimental Setup}
\label{ch:methodology}

% This chapter explains every model and every shared experimental choice
% (preprocessing, hyperparameters, evaluation protocol) that applies across
% all three datasets. Dataset-specific preprocessing details that don't
% generalize (e.g. exact feature counts, window lengths) can either stay
% here at a high level or move into the per-dataset Results subsections --
% decide per case once the prose is drafted.

\section{Anomaly Detection Paradigms}
% Short framing paragraph: novelty detection (train on benign only) vs.
% outlier detection, and the four paradigms under comparison --
% boundary-based, partition-based, reconstruction-based, temporal (predictive
% and reconstruction).

\subsection{One-Class SVM}
% Hypersphere/boundary intuition, nu-SVM formulation, kernel choice.

\subsection{Isolation Forest}
% Partition-based isolation, path length as anomaly score.

\subsection{Feedforward Autoencoder}
% Reconstruction-based, bottleneck compression, MSE score.

\subsection{LSTM Autoencoder}
% Sequence-to-vector-to-sequence reconstruction; contrast with predictive LSTM below.

\subsection{Predictive LSTM}
% Next-step prediction objective; prediction residual as anomaly score.
% Note explicitly how this differs in objective (not just architecture)
% from the LSTM Autoencoder above -- this distinction is a key result later.

\subsection{Transformer-Based Autoencoders}
% Two variants used: bottleneck transformer (global-average-pool + squeeze/expand
% bottleneck) and causal transformer (masked self-attention, predictive objective).
% Mention the originally-planned masked/token-reconstruction transformer
% (IMPLEMENTATION_PLAN.md \S2.5) and its current status (not yet run on Yahoo).

\section{Datasets}
% One short paragraph per dataset here (name, size, source, why chosen),
% deferring full preprocessing detail to \S\ref{sec:preprocessing} and
% deferring results/discussion to Chapter~\ref{ch:results}.

\subsection{Credit Card Fraud Detection (Dataset A)}
\subsection{Yahoo Webscope S5 (Dataset B)}
\subsection{CIC-IDS2017 (Dataset C)}

\section{Preprocessing}
\label{sec:preprocessing}
% Shared principles first (train-on-benign-only split, scaler fit on train
% only, no leakage), then one subsection per dataset for the concrete pipeline.

\subsection{Tabular Preprocessing (Credit Card)}
\subsection{Time-Series Windowing (Yahoo S5)}
\subsection{Network Flow Preprocessing (CIC-IDS2017)}
% Include the flat / temporal / host-context variants and why context
% features were introduced.

\section{Hyperparameter Choices and Architecture Standardization}
% Pull directly from mds/architecture_diffs.md: table per model (OC-SVM,
% IForest, MLP AE, LSTM, LSTM AE, Bottleneck Transformer, Causal Transformer)
% showing the values used per dataset, plus the justification for why values
% differ (e.g. bottleneck ratio kept constant, not absolute width).
% State explicitly: no hyperparameter sweeping was performed by design --
% goal is paradigm comparison under principled, fixed configurations, not
% per-model optimality.

\section{Evaluation Protocol}
% Metrics: PR-AUC (Average Precision) as primary metric, ROC-AUC for
% comparability, F2-maximizing threshold via sweep\_threshold() as secondary/
% operational metric. Justify PR-AUC choice under class imbalance.
% State the "no Point Adjustment" methodological choice and cite Kim et al.
% (2022) critique of PA-adjusted F1.
% State the oracle-threshold limitation (threshold selected from the test
% set's own PR curve) and why it's acceptable for a comparison-only study.
% Mention the trivial baseline (\|w\|\_2) sanity check used on Yahoo.
